% Part: incompleteness
% Chapter: theories-computability
% Section: introduction

\documentclass[../../../include/open-logic-section]{subfiles}

\begin{document}

\olfileid[tr]{inc}{tcp}{int}

\olsection{Giriş}

\begin{editorial}
Bu bölüm yeniden yazılmalıdır.
\end{editorial}

Elimizde şunlar var:
\begin{enumerate}
\item Bir fonksiyonun $\Th{Q}$ içinde temsil edilebilir olmasının ne anlama
  geldiğinin tanımı (\olref[req][int]{defn:representable-fn})
\item bir bağıntının $\Th{Q}$ içinde temsil edilebilir olmasının ne anlama
  geldiğinin tanımı (\olref[req][rel]{defn:representing-relations})
\item $\Th{Q}$ içinde temsil edilebilen fonksiyonların tam olarak
  hesaplanabilir fonksiyonlar olduğunu söyleyen bir teorem
  (\olref[req][int]{thm:representable-iff-comp})
\item $\Th{Q}$ içinde temsil edilebilen bağıntıların tam olarak
  hesaplanabilir bağıntılar olduğunu söyleyen bir teorem
  (\olref[req][rel]{thm:representing-rels})
\end{enumerate}

Bir {\em kuram}, türetim altında kapalı bir tümceler kümesidir; başka bir
deyişle $T$, $!A$'yı ispatladığında $!A$'nın $T$ içinde bulunması özelliğine
sahiptir. Bir kuramı, bir tümceler topluluğu ile bu tümcelerin gerektirdiği
her şeyin birlikte oluşturduğu bir bütün olarak düşünmek herhâlde en
iyisidir. Bundan böyle $\Th{Q}$ ile, \olref[req][int]{sec}'deki sekiz
aksiyomdan türetilebilen tümceler kümesinden oluşan {\em kuramı}
kastediyoruz. $\Th{Q}$ formüllerini sayılarla kodlayabildiğimizi anımsayın;
$!A$ böyle bir formülse $\Gn{!A}$, $!A$'yı kodlayan sayı olsun. Artık bu
kodlamaya göre çeşitli formül kümelerinin hesaplanabilir olup olmadığını
sorabiliriz.

\end{document}
